%% sections/04-spiked-covariance-bbp.tex
%% Section 4: Spiked Covariance Matrices (~11 pages)
%% Source notes:
%%   - v5 LaTeX notes: spiked covariance, BBP, simulations, overlap, clustering.
%%   - Handwritten Gustavo Notes (April 17, April 27, April 30, May 14).
%%   - Yao, Zheng, Bai (2015) Chapter 11: spiked sample covariance matrices.
%%   - F\'eral and P\'ech\'e (2009) Theorem 1.6 for fluctuations.
%%   - Existing figures in LaTeX notes v5/figures/:
%%       smoothed_m1_strong_vs_weak.png, m1_strong_spike.png, m1_weak_spike.png,
%%       eigenvector_overlap_distribution_n1200.png,
%%       m2_two_strong_spikes.png, m2_mixed_spikes.png.

\section{Spiked Covariance Matrices and the BBP Transition}\label{sec:spiked}

\todo{Opening paragraph: the Marchenko-Pastur law of
Section~\ref{subsec:scm-highdim} described the null case \(\Sigma = I_p\),
where all population eigenvalues are equal. In applications the population
covariance carries structure: a small number of distinguished directions
of high variance riding above a flat noise baseline. The
\emph{spiked covariance model} of \citet{Johnstone2001} formalizes this,
\[
  \Sigma = \diag(\alpha_1, \dots, \alpha_m, 1, \dots, 1),
\]
with \(m\) fixed as \(p \to \infty\). The central question is when these
population spikes are detectable from the sample spectrum, and when the
sample eigenvectors carry meaningful information about the population
spike directions. The answer is a clean phase transition.}

\subsection{Johnstone's Spiked Model and the Spike-Location Map}\label{subsec:johnstone}

\todo{Content to fill in (from v5 + Yao-Zheng-Bai Ch.\ 11 + Gustavo notes):
- Definition: Johnstone's spiked model
  \(\Sigma = \diag(\alpha_1, \dots, \alpha_m, 1, \dots, 1)\) with
  \(H = \delta_1\) for the bulk limiting distribution.
- Definition: spike-location map
  \(\Psi(\alpha) = \alpha + \frac{y \alpha}{\alpha - 1}\).
- Compute \(\Psi'(\alpha) = 1 - \frac{y}{(\alpha - 1)^2}\).
- Lemma: \(\Psi'(\alpha) > 0 \iff |\alpha - 1| > \sqrt{y}\); for large spikes
  \(\alpha > 1\) this becomes \(\alpha > 1 + \sqrt{y}\).
- Verify boundary behavior: \(\Psi(1 \pm \sqrt{y}) = (1 \pm \sqrt{y})^2\),
  so the spike map agrees with the MP edges at the threshold.}

\subsection{The BBP Phase Transition for Eigenvalues}\label{subsec:bbp}

\todo{Content to fill in:
- Theorem (Yao-Zheng-Bai Corollary 11.4 specialized to Johnstone):
  For a large spike \(\alpha_j > 1 + \sqrt{y}\),
  \[
    \lambda_i(\Sigmahat_n) \convas \Psi(\alpha_j), \quad i \in J_j.
  \]
  For \(1 < \alpha_j \le 1 + \sqrt{y}\),
  \[
    \lambda_i(\Sigmahat_n) \convas (1 + \sqrt{y})^2.
  \]
  Brief sketch (cite \citet{BaikSilverstein2006} for the proof).
- Worked example: \(y = 1/2\), so \(1 + \sqrt{y} \approx 1.707\) and edge
  \((1 + \sqrt{y})^2 \approx 2.914\). For \(\alpha = 2.5\),
  \(\Psi(2.5) \approx 3.333\). For \(\alpha = 1.4\), \(\alpha < 1.707\) so
  the spike is absorbed into the edge.
- Simulations (existing figures from v5):
  Figure: smoothed m=1 simulation, strong vs.\ weak spike
  (\texttt{smoothed\_m1\_strong\_vs\_weak.png}).
  Figure: raw top eigenvalue distributions
  (\texttt{m1\_strong\_spike.png}, \texttt{m1\_weak\_spike.png}).
  Figure (optional): \(m=2\) cases
  (\texttt{m2\_two\_strong\_spikes.png}, \texttt{m2\_mixed\_spikes.png}).
Cite \citet{BaikBenArousPeche2005, YaoZhengBai2015}.}

\subsection{Fluctuations of Spiked Eigenvalues}\label{subsec:fluctuations}

\todo{Content to fill in (from Gustavo's promised sketch + v5):
- Statement (fundamental spike, qualitative): for
  \(\alpha > 1 + \sqrt{y}\) in the Gaussian Johnstone model,
  \[
    \sqrt{n}(\lambda_{\max}(\Sigmahat_n) - \Psi(\alpha)) \convd
    \Normal{0, \sigma_{\alpha, y}^2}.
  \]
  Cite \citet[Thm.~11.11]{YaoZhengBai2015} (fundamental spikes) and
  \citet[Thm.~1.6]{FeralPeche2009} (cleanest diagonal case).
- Heuristic: a separated eigenvalue is a smooth functional of \(\Sigmahat_n\),
  so a delta-method argument gives \(\sqrt{n}\) Gaussian fluctuations
  (compare with Section~\ref{subsec:scm-lowdim}(e)).
- Non-fundamental spike: top eigenvalue attaches to the edge,
  \(\lambda_{\max} - \lambda_+ = O_p(p^{-2/3})\), and the fluctuation is
  expected to be Tracy-Widom type (cite \citet{Johnstone2001}).
- Honest note: the critical case \(\alpha = 1 + \sqrt{y}\) is not resolved
  in our current treatment; cite open literature.
- Simulation: histogram and Gaussian Q-Q for strong-spike fluctuations
  (notebook 06).}

\subsection{Eigenvector Alignment Beyond the Threshold}\label{subsec:eigenvector}

\todo{Content to fill in (from v5 + Yao-Zheng-Bai Thm 11.5/Cor 11.6):
- Set-up: in the one-spike model, the population spike direction is
  \(\vect e_1\). Let \(\widehat{\vect u}_1\) be the top sample eigenvector.
  The squared overlap \(|\inner{\widehat{\vect u}_1}{\vect e_1}|^2 =
  \cos^2 \theta\) removes the sign ambiguity.
- Qualitative theorem (Yao-Zheng-Bai Thm.~11.5, Cor.~11.6):
  strong spike \(\Rightarrow\) overlap converges to a positive constant;
  weak spike \(\Rightarrow\) overlap converges to zero.
- The one-spike Johnstone overlap formula:
  \[
    \rho(\alpha, y) = \frac{1 - \dfrac{y}{(\alpha - 1)^2}}
                           {1 + \dfrac{y}{\alpha - 1}}.
  \]
  Positive precisely when \(\alpha > 1 + \sqrt{y}\).
- Simulation (figure already in v5/figures/):
  \texttt{eigenvector\_overlap\_distribution\_n1200.png} shows strong-spike
  overlaps clustered around 0.55--0.65 and weak-spike overlaps near 0.
  Include and discuss.}

\subsection{Sketch of the Connection to Spectral Clustering}\label{subsec:clustering}

\todo{Content to fill in (from Gustavo's email + v5 + von Luxburg 2007):
- Setup: data points \(\vect x_1, \dots, \vect x_n \in \R^p\), kernel
  similarity \(K_{ij} = f(\norm{\vect x_i - \vect x_j}^2 / p)\).
- Definition: degree matrix \(D = \diag(K \vect 1)\) and normalized
  Laplacian-style kernel \(L = n D^{-1/2} K D^{-1/2}\).
- Heuristic decomposition:
  \(L = \text{noise bulk} + \text{low-rank class signal}\),
  in direct analogy with \(\Sigmahat_n = \text{noise} + \text{spike}\).
- Conclusion: class structure is recoverable from \(L\) precisely when the
  signal-to-noise ratio crosses an analogue of the BBP threshold; the
  informative eigenvectors of \(L\) then provide the clustering. Cite
  \citet{vonLuxburg2007} for the algorithmic side and
  \citet{CouilletBenaych2016} for the high-dimensional analysis.
- Toy simulation: notebook 08, a two-class Gaussian mixture and the
  top non-trivial eigenvector of \(L\).}

\clearpage
